\subsubsection{Sobolev Spaces and Fourier Transform}\label{subsec:sobolev-fourier-characterization}

We now consider Sobolev spaces of all of $\mathbb{R}^m$. Clearly, it follows from Theorem~\ref{thm:fourier-inner-product-preservation} that the Fourier transform maps $L^2(\mathbb{R}^m)$ to itself; indeed, it is an isometry in $L^2(\mathbb{R}^m)$. Moreover, the Fourier transform of $D^\alpha u$ is $(i\bm{\xi})^\alpha\hat u$. Hence we immediately obtain the following result.

\begin{Theorem}
\label{thm:sobolev-fourier-weighted-homeomorphism}
The Fourier transform $\mathcal{F}$ is a homeomorphism from $H^k(\mathbb{R}^m)$ onto the weighted space $L_w^2(\mathbb{R}^m)$ (cf. Definition~\ref{def:weighted-lp-two-space}), where $w(\bm{\xi})=1+|\bm{\xi}|^{2k}$.
\end{Theorem}

We shall use the notation $L_k^2$ to denote this weighted $L^2$-space. It is easy to see that $\mathcal{S}(\mathbb{R}^m)$ is dense in $L_k^2(\mathbb{R}^m)$. Theorem~\ref{thm:sobolev-fourier-weighted-homeomorphism} then implies that $\mathcal{S}(\mathbb{R}^m)$, and hence also $\mathcal{D}(\mathbb{R}^m)$, is dense in $H^k(\mathbb{R}^m)$.

\begin{Corollary}
\label{cor:full-space-test-functions-sobolev-density}
$\mathcal{D}(\mathbb{R}^m)$ is dense in $H^k(\mathbb{R}^m)$.
\end{Corollary}

Another application of the theorem is the definition of fractional order Sobolev spaces.

\begin{Definition}
\label{def:fractional-order-sobolev-fourier}
We say that $u\in H^s(\mathbb{R}^m)$, $s\in\mathbb{R}^+$, if $\mathcal{F}[u]$ is in the weighted $L^2$-space $L_w^2(\mathbb{R}^m)=:L_s^2(\mathbb{R}^m)$ with $w(\bm{\xi})=1+|\bm{\xi}|^{2s}$.
\end{Definition}

There is an intrinsic characterization of the fractional Sobolev spaces, which is basically an $L^2$-analogue of Hölder continuity. It can be shown that an equivalent inner product on $H^s(\mathbb{R}^m)$ is given by
\begin{equation}\label{eq:fractional-sobolev-intrinsic-inner-product}
\begin{aligned}
(u,v)_s&=(u,v)_{[s]}\\
&\quad+\sum_{|\alpha|=[s]}\int_{\mathbb{R}^m}\int_{\mathbb{R}^m}
\dfrac{(D^\alpha\bar u(\bm{x})-D^\alpha\bar u(\bm{y}))(D^\alpha v(\bm{x})-D^\alpha v(\bm{y}))}{|\bm{x}-\bm{y}|^{m+2(s-[s])}}\,\mathrm{d}\bm{x}\,\mathrm{d}\bm{y}.
\end{aligned}
\end{equation}
Here $[s]$ denotes the largest integer $\leq s$, and the integral in \eqref{eq:fractional-sobolev-intrinsic-inner-product} is to be interpreted, as usual, as a limit for sequences of smooth functions which approximate $u$ and $v$, respectively. Clearly, a definition in terms of \eqref{eq:fractional-sobolev-intrinsic-inner-product} can be extended immediately to arbitrary $\Omega$.
